\subsection{2025}

\subsubsection{May 29}

\phantomsection
\addcontentsline{toc}{paragraph}{Quantum Circuits}
\eexercise[Quantum Circuits][7]
\begin{itemize}
    \item Consider the following quantum circuit defined on 2 qubits. Compute the probability of measuring each possible state at the end of the circuit.
        \begin{center}
            \begin{quantikz}
                \lstick{$q_1 = \ket{1}$} & \gate{H} & \ctrl{1} & \targ{}   & \gate{Z} & \meter{} \\
                \lstick{$q_2 = \ket{0}$} & \gate{X} & \gate{S} & \ctrl{-1} & \gate{Y} & \meter{}
            \end{quantikz}
        \end{center}

    \item Describe what phase kickback is and when it happens. State if phase kickback occurs in the above quantum circuit and, if yes, where it occurs and why.
\end{itemize}
Motivate your answer by showing all stages of the computation.

\addcontentsline{toc}{paragraph}{Measurement}
\highspace
\esolution[Measurement] First, we divide the circuit into stages and compute the state of the system after each stage:
\begin{center}
    \begin{quantikz}[
            slice all,
            slice titles={\ket{\psi_{\col}}},
            column sep=1cm,
        ]
        \lstick{$q_1 = \ket{1}$}\slice{\ket{\psi_{0}}} & \gate{H} & \ctrl{1} & \targ{}   & \gate{Z} & \meter{} \\
        \lstick{$q_2 = \ket{0}$}                       & \gate{X} & \gate{S} & \ctrl{-1} & \gate{Y} & \meter{}
    \end{quantikz}
\end{center}
\noindent
Let's compute the state of the system after each stage:
\begin{itemize}
    \item The initial state of the system is:
        \begin{equation*}
            \boxed{\ket{\psi_{0}} = \ket{1} \otimes \ket{0} = \ket{10}}
        \end{equation*}

    \item Apply the Hadamard gate to the first qubit and the $X$ gate to the second qubit:
        \begin{equation*}
            \boxed{\ket{\psi_{1}} = (H \otimes X) \ket{10} = \frac{1}{\sqrt{2}} \left(\ket{0} - \ket{1}\right) \otimes \ket{1} = \frac{1}{\sqrt{2}} \left(\ket{01} - \ket{11}\right)}
        \end{equation*}

    \item Apply the controlled-$S$ gate to the second qubit. Remember that the $S$ gate is defined as:
        \begin{equation*}
            S =
            \begin{pmatrix}
                1 & 0 \\
                0 & i
            \end{pmatrix}
        \end{equation*}
        The controlled-$S$ gate means that if the first qubit is in state $\ket{1}$, we apply the $S$ gate to the second qubit. Therefore, we have:
        \begin{itemize}
            \item For \ket{01}, the first qubit is in state $\ket{0}$, so we do not apply the $S$ gate to the second qubit.

            \item For \ket{11}, the first qubit is in state $\ket{1}$, so we apply the $S$ gate to the second qubit. The second qubit is in state $\ket{1}$, so we have:
                \begin{equation*}
                    S\ket{1} = i\ket{1}
                \end{equation*}
        \end{itemize}
        Combining these results, we have:
        \begin{equation*}
            \boxed{\ket{\psi_{2}} = \frac{1}{\sqrt{2}} \left(\ket{01} - i\ket{11}\right)}
        \end{equation*}

    \item Apply the CNOT gate, where the control is the second qubit and the target is the first qubit. The CNOT gate applies the Pauli-$X$ gate to the target qubit if the control qubit is in state $\ket{1}$. Since both terms in $\ket{\psi_{2}}$ have the second qubit in state $\ket{1}$, we apply the $X$ gate to the first qubit in both terms:
        \begin{itemize}
            \item For \ket{01}, the calculation is simple:
                \begin{equation*}
                    X\ket{0} = \ket{1}
                \end{equation*}
                So, the first term becomes \ket{11}.

            \item For \ket{11}, we should avoid the mistake of assuming that the $i$ factor influences the CNOT gate's result. Instead, we apply the $X$ gate to the first qubit.
                \begin{equation*}
                    X\ket{1} = \ket{0}
                \end{equation*}
                So, the second term becomes $-i\ket{01}$.
        \end{itemize}
        Combining these results, we have:
        \begin{equation*}
            \boxed{\ket{\psi_{3}} = \frac{1}{\sqrt{2}} \left(\ket{11} - i\ket{01}\right)}
        \end{equation*}

    \item Before measurement, we apply the $Z$ gate to the first qubit and the $Y$ gate to the second qubit. The $Z$ gate is defined as:
        \begin{equation*}
            Z =
            \begin{pmatrix}
                1 & 0 \\
                0 & -1
            \end{pmatrix}
        \end{equation*}
        And the $Y$ gate is defined as:
        \begin{equation*}
            Y =
            \begin{pmatrix}
                0 & -i \\
                i & 0
            \end{pmatrix}
        \end{equation*}
        The $Z$ gate flips the sign only for the \ket{11} state:
        \begin{equation*}
            Z\ket{11} = -\ket{11}
        \end{equation*}
        And the $Y$ gate (1) swaps the states, (2) flips the sign of the \ket{0}, and (3) adds a factor of $i$:
        \begin{align*}
            Y(-\ket{11}) &\xrightarrow{1} (-\ket{10}) \xrightarrow{2} \ket{10} \xrightarrow{3} i\ket{10} \\
            Y(-i\ket{01}) &\xrightarrow{1} (-i\ket{00}) \xrightarrow{2} (i\ket{00}) \xrightarrow{3} (\underbrace{i \cdot i}_{-1}\ket{00}) = -\ket{00}
        \end{align*}
        Combining these results, we have:
        \begin{equation*}
            \boxed{\ket{\psi_{4}} = \frac{1}{\sqrt{2}} \left(i\ket{10} - \ket{00}\right)}
        \end{equation*}

    \item Finally, we measure the system. To simplify the understanding of the measurement, we can expand the state $\ket{\psi_{4}}$ as follows:
        \begin{equation*}
            \ket{\psi_{4}} = -\frac{1}{\sqrt{2}}\ket{00} + 0\ket{01} + \frac{i}{\sqrt{2}} \ket{10} + 0\ket{11}
        \end{equation*}
        The probability of measuring each possible state is given by the square of the absolute value of the coefficients:
        \begin{itemize}
            \item Probability of measuring \ket{00}:
                \begin{equation*}
                    P(\ket{00}) = \left|-\frac{1}{\sqrt{2}}\right|^2 = \frac{1}{2}
                \end{equation*}
            \item Probability of measuring \ket{01}:
                \begin{equation*}
                    P(\ket{01}) = \left|0\right|^2 = 0
                \end{equation*}
            \item Probability of measuring \ket{10}:
                \begin{equation*}
                    P(\ket{10}) = \left|\frac{i}{\sqrt{2}}\right|^2 = \underbrace{\left(
                            -\frac{i}{\sqrt{2}} \cdot \frac{i}{\sqrt{2}}
                    \right)}_{|z|^2 = \bar{z} \cdot z} = -\frac{i^2}{2} = \frac{1}{2}
                \end{equation*}
            \item Probability of measuring \ket{11}:
                \begin{equation*}
                    P(\ket{11}) = \left|0\right|^2 = 0
                \end{equation*}
        \end{itemize}
\end{itemize}
Thus we have the \hl{$50$\% probability of measuring either $\ket{00}$ or $\ket{10}$}, and $0$\% probability of measuring \ket{01} or \ket{11}.

\highspace
\esolution*[Phase Kickback]
Phase kickback is a quantum phenomenon that can occur in a controlled operation. In particular, consider a controlled-$U$ gate where the target qubit is in an eigenstate $\ket{v}$ of $U$:

\begin{equation*}
    U\ket{v}=e^{i\phi}\ket{v}
\end{equation*}
If the control qubit is in a superposition, for example:
\begin{equation*}
    \frac{\ket{0}+\ket{1}}{\sqrt{2}}
\end{equation*}
Then:
\begin{equation*}
    CU
    \left(
        \frac{\ket{0}+\ket{1}}{\sqrt{2}}
        \otimes \ket{v}
    \right)
    =
    \frac{\ket{0}+e^{i\phi}\ket{1}}{\sqrt{2}}
    \otimes \ket{v}
\end{equation*}
Therefore, the target qubit remains in the same eigenstate $\ket{v}$, while the phase $e^{i\phi}$ appears on the $\ket{1}$ component of the control qubit. This is called \textbf{phase kickback}.

\highspace
In the given circuit, immediately before the controlled-$S$ gate we have:
\begin{equation*}
    \ket{\psi_1}
    =
    \frac{\ket{0}-\ket{1}}{\sqrt{2}}
    \otimes \ket{1}
\end{equation*}
The target qubit is $\ket{1}$, which is an eigenstate of the $S$ gate because:
\begin{equation*}
    S\ket{1}=i\ket{1}
\end{equation*}
Hence:
\begin{equation*}
    CS
    \left(
        \frac{\ket{0}-\ket{1}}{\sqrt{2}}
        \otimes \ket{1}
    \right)
    =
    \frac{\ket{0}-i\ket{1}}{\sqrt{2}}
    \otimes \ket{1}
\end{equation*}
Thus, phase kickback occurs at the controlled-$S$ gate: the target remains in state $\ket{1}$, while the eigenvalue $i$ of $S$ modifies the relative phase of the control qubit.